% titel: Blatt 4, A5: Photoelektrischer Effekt an Molybdän
% Teil II: der durchgerechnete Weg.

\subsection*{Gegeben und gesucht}
Grenzfrequenz Mo: $\nu_0 = \qty{1.09e15}{s^{-1}}$.
a) Mindestenergie (Austrittsarbeit $W$),
b) zugehörige Wellenlänge $\lambda_0$,
c) maximale $E_{\text{kin}}$ bei Bestrahlung mit $\lambda = \qty{120}{nm}$.

\subsection*{Ansatz}
$W = h\nu_0$; $\lambda_0 = c/\nu_0$; $E_{\text{kin}} = hc/\lambda - W$.

\subsection*{Rechnung}
\begin{align*}
  W &= 6.626\times10^{-34}\times 1.09\times10^{15}
  = 7.22\times10^{-19}\ \mathrm{J}
  = 4.51\ \mathrm{eV}, \\[2pt]
  \lambda_0 &= \frac{2.998\times10^{8}}{1.09\times10^{15}}
  = 2.75\times10^{-7}\ \mathrm{m}
  = \qty{275}{nm}, \\[2pt]
  E_{\text{photon}} &= \frac{h\,c}{\lambda}
  = \frac{1.987\times10^{-25}}{1.20\times10^{-7}}
  = 1.66\times10^{-18}\ \mathrm{J}
  = 10.33\ \mathrm{eV}, \\[2pt]
  E_{\text{kin}} &= E_{\text{photon}} - W
  = 1.66\times10^{-18} - 7.22\times10^{-19}
  = 9.33\times10^{-19}\ \mathrm{J}
  = 5.82\ \mathrm{eV}.
\end{align*}

\begin{ergebnis}
a) Austrittsarbeit $W \approx \qty{7.2e-19}{J} = \qty{4.5}{eV}$;
b) $\lambda_0 \approx \qty{275}{nm}$;
c) $E_{\mathrm{kin,max}} \approx \qty{9.3e-19}{J} = \qty{5.8}{eV}$.
\end{ergebnis}

\begin{kontrolle}
Die Grenzfrequenz $1.09\times10^{15}$ s$^{-1}$ entspricht einer
Abschneidewellenlänge von 275 nm (UV); kürzeres Licht (120 nm) hat mehr
Energie, der Überschuss $10.3 - 4.5 = 5.8\ \mathrm{eV}$ erscheint als
kinetische Energie des Elektrons. Plausibel und größenordnungsmäßig typisch
für den Photoeffekt.
\end{kontrolle}

\begin{merksatz}
Photoeffekt: W = h$\nu_0$; $E_{\text{kin}} = hc/\lambda - W$. Unterhalb der
Grenzfrequenz (oberhalb $\lambda_0$) gibt es keine Emission.
\end{merksatz}
